% ============================================================
%  PipedPeer -- journal manuscript (Elsevier elsarticle)
%  Target: Journal of Parallel and Distributed Computing.
%  Every number traces to docs/benchmarks/ or to the commit named in a
%  "% source:" comment beside it. Remove those comments before an
%  arXiv or source upload: they are notes for the authors.
%  Build: tectonic manuscript.tex   (or pdflatex twice)
% ============================================================
\documentclass[5p,times,twocolumn]{elsarticle}

\usepackage{amsmath,amssymb}
\usepackage{graphicx}
\usepackage{booktabs}
\usepackage{tabularx}
\usepackage{url}
\usepackage{algorithm}
\usepackage{algpseudocode}
\usepackage{xcolor}
\usepackage[hidelinks]{hyperref}

\newcommand{\figpath}[1]{../../diagrams/#1}
% Marks a field only the authors can fill. The manuscript must contain none
% of these when it is submitted.
\newcommand{\authorfill}[1]{\textcolor{red}{\textbf{[AUTHORS: #1]}}}
\newcolumntype{L}[1]{>{\raggedright\arraybackslash}p{#1}}
\newcolumntype{C}[1]{>{\centering\arraybackslash}p{#1}}
\renewcommand{\tabularxcolumn}[1]{>{\raggedright\arraybackslash}p{#1}}

% Long code words (multiprocessing, concurrent.futures) would otherwise
% overrun the margin.
\emergencystretch=3em

\journal{Journal of Parallel and Distributed Computing}

\begin{document}

\begin{frontmatter}

\title{PipedPeer: Running Unmodified Python Across a Small Group's Machines
Without a Configured Cluster}

\author[spit]{Yateen Vaviya\corref{cor1}\fnref{eq}}
\ead{yateen.vaviya23@spit.ac.in}
\author[spit]{Vinayak Yadav\fnref{eq}}
\ead{vinayak.yadav23@spit.ac.in}
\author[spit]{Sahil Gupta\fnref{eq}}
\ead{sahil.gupta23@spit.ac.in}
\author[spit]{Rupali Sawant}
\ead{rupali_sawant@spit.ac.in}
\cortext[cor1]{Corresponding author.}
\fntext[eq]{These authors contributed equally to this work.}

\affiliation[spit]{organization={Department of Computer Engineering, Bharatiya Vidya Bhavan's Sardar Patel Institute of Technology},
  addressline={Bhavan's Campus, Old D N Nagar, Munshi Nagar, Andheri West},
  city={Mumbai},
  postcode={400058},
  state={Maharashtra},
  country={India}}

\begin{abstract}
Cluster frameworks (Spark, Ray, Dask) assume a configured cluster, a
program written against their interface, and machines that can reach one
another. This leaves a small group's idle laptops, desktops and phones out of reach,
especially behind different home or campus routers. PipedPeer runs an unmodified Python script on other machines.
Every machine runs the same daemon and there is no head node. The script's
environment is built as a Nix closure pinned to one package-set revision, and a
peer receives only the store paths it lacks. Parallel maps, joblib, pandas
grouping and joins, NumPy products and PyTorch data-parallel training are
substituted at import time. A map is divided by a makespan model fitted to
measured peer rates, and numerical work moves only when a cost model predicts a
gain. Machines on different networks meet through an introducer that every
daemon serves and then connect over direct, authenticated QUIC links; no job
data is relayed. We evaluate PipedPeer in container labs, on two laptops in
different homes with a phone, and on a second pair of machines. A repeated
submission whose closure the worker holds takes 0.83\,s against 23.4\,s for the first. A
mixed CPU script took 13.3--14.6\,s on two laptops against 17.4--18.1\,s on
one. Periodic weight averaging every 32 steps kept validation loss within 1\%
of one device on three test models. Small training jobs were not faster on
two machines than on one GPU. The evaluation covers two or three physical
machines, not a controlled scaling study.
\end{abstract}

\begin{keyword}
distributed computing \sep peer-to-peer systems \sep NAT traversal \sep
transparent interception \sep reproducible environments \sep
heterogeneous scheduling \sep data-parallel training
\end{keyword}

\end{frontmatter}


% ============================================================
\section{Introduction}
\label{sec:intro}

The gap between running \texttt{python script.py} on one machine and running
the same work on several is large in operational effort and often small in the
computation involved. Apache Spark~\cite{zaharia2010spark},
Ray~\cite{moritz2018ray} and Dask~\cite{rocklin2015dask} presuppose a configured
cluster with a head node, matching software on every worker, a scheduler, and a
network on which every machine reaches every other. Remote execution over
\texttt{ssh} avoids the setup but offers no environment management, placement,
isolation or fault tolerance.

For a small group --- a student team, a laboratory, a household --- three
obstacles stand between a script and the machines around it, and they are
usually treated separately:
\begin{enumerate}
\item \emph{Environment}: a remote machine must be made to hold the same
interpreter and libraries as the local one.
\item \emph{Code}: using more than one machine normally means rewriting the
program against a framework's API.
\item \emph{Reachability}: machines behind different home or campus routers
accept no inbound connections, and the usual remedies --- a VPN, port
forwarding or a relay --- are either configuration or a bottleneck.
\end{enumerate}

PipedPeer addresses all three for Python on Linux. A user runs
\texttt{pipedpeer run script.py}; the client builds the script's environment as
a Nix closure~\cite{dolstra2004nix}, picks a machine, ships the missing store
paths and the workspace, and runs the job in a sandbox. A Python package
imported before the script's first line substitutes the parallel primitives the
script already uses, and spreads their work across machines when a model of
the cluster predicts that doing so is faster. Machines on different networks
join with one invite code and connect directly.

\paragraph{Contributions}
\begin{enumerate}
\item \textbf{Peer-to-peer closure distribution} (Section~\ref{sec:env}):
signed closures pushed between peers, diffed by store path, with distributed
work offered only to peers that can run the closure, and environments of a
foreign architecture built by the receiver.
\item \textbf{Import-time interception with measured work sharing}
(Sections~\ref{sec:intercept}--\ref{sec:cost}): parallel maps divided by a
makespan model over measured rates, with speculative copies of late parts, and
numerical work gated by a cost inequality.
\item \textbf{Joining across networks without a relay} (Section~\ref{sec:nat}):
an introducer embedded in every daemon that only exchanges addresses, router
port mapping, a candidate race on one shared socket, and a birthday-collision
fallback.
\item \textbf{Placement without a scheduler process}
(Section~\ref{sec:placement}): filters and a bounded score computed in the
submitting client, with a three-phase lease.
\item \textbf{Speed-weighted data-parallel training through the daemons}
(Section~\ref{sec:train}), with exact or periodic-averaging synchronisation
and an interval chosen per ring size from accuracy measurements.
\item \textbf{An evaluation that keeps emulated, laboratory and physical-machine
results apart} (Section~\ref{sec:eval}), including the cases in which
distribution was slower than one machine.
\end{enumerate}

\paragraph{Scope} PipedPeer targets two to a handful of machines owned by one
group that trusts itself. It is not a replacement for the frameworks above at
scale, and the evaluation does not establish speedup beyond three machines
(Section~\ref{sec:threats}).

% ============================================================
\section{Background and related work}
\label{sec:related}

\paragraph{Cluster frameworks} Spark~\cite{zaharia2010spark},
Ray~\cite{moritz2018ray} and Dask~\cite{rocklin2015dask} are mature and far
more capable at scale. Ray ships a drop-in \texttt{multiprocessing.Pool} and a
\texttt{joblib} backend, and Modin~\cite{petersohn2020modin} provides
distributed pandas behind the pandas API; each needs a changed import and a
running cluster. Kubernetes~\cite{burns2016borg} and HTCondor~\cite{litzkow1988condor}
schedule whole jobs through a central manager and explicit job descriptions.

\paragraph{Drop-in parallelism for Python} Several systems keep the
\texttt{multiprocessing} interface and change only where the work runs.
Fiber~\cite{zhi2020fiber} replaces \texttt{multiprocessing} with processes on a
cluster manager; PyWren~\cite{jonas2017pywren} and its successor
Lithops~\cite{sampe2021lithops} run maps on serverless functions, and Lithops
exposes a \texttt{multiprocessing}-compatible API; Arjona
et~al.~\cite{arjona2023transparent} extend it to run unmodified
\texttt{multiprocessing} applications, shared state included, on serverless
functions. Parsl~\cite{babuji2019parsl}
and funcX~\cite{chard2020funcx} turn decorated Python functions into tasks on
remote resources. All of these require a provisioned back end (a cluster manager or a cloud
account), and all but Arjona et~al.'s an import or decorator change. PipedPeer's
substitution happens before the script's first statement, so no line changes,
and it decides per call whether to distribute at all. The closest system in
intent is therefore Arjona et~al.'s; the differences are that PipedPeer
needs no back end beyond the peers themselves, ships the environment to them,
crosses NATs between them, and may decline to distribute.

\paragraph{Reproducible environments} Nix~\cite{dolstra2004nix,dolstra2006thesis}
and Guix~\cite{courtes2013guix} name each build by a hash of its inputs, which
makes a store path a stable name for an environment; Spack~\cite{gamblin2015spack}
targets HPC software stacks; container images~\cite{merkel2014docker,kurtzer2017singularity}
ship a whole filesystem. PipedPeer uses Nix store paths and Nix's own binary
cache format, and adds peer-to-peer pushing, path-level diffs between peers,
and a rule that only peers able to run the closure receive distributed work.

\paragraph{Harvesting idle machines} Condor~\cite{litzkow1988condor} scavenged
idle workstations; SETI@home~\cite{anderson2002seti} and
BOINC~\cite{anderson2004boinc} harvest volunteers' machines for specific
projects. Smartphones have been studied both as offload clients ---
MAUI~\cite{cuervo2010maui} and CloneCloud~\cite{chun2011clonecloud} move work
\emph{off} a phone --- and as compute donors: CWC~\cite{arslan2012cwc} schedules
jobs on phones while they charge, and Honeybee~\cite{fernando2012honeybee}
distributes work among nearby phones. PipedPeer's phone support is narrower
than either line: an Android phone under Termux takes parallel-map items
(Section~\ref{sec:phone}).

\paragraph{NAT traversal} Reflexive addresses are discovered as in
STUN~\cite{rfc8489}; port mappings are requested with PCP~\cite{rfc6887},
NAT-PMP~\cite{rfc6886} and UPnP-IGD~\cite{upnpigd}; ICE~\cite{rfc8445}
standardises candidate gathering and pairing; simultaneous UDP probing from the
socket that will carry the connection is the hole punching of Ford
et~al.~\cite{ford2005p2p}; and the birthday-collision technique for a
port-randomising NAT follows Anderson~\cite{anderson2020nat}. Links are
QUIC~\cite{rfc9000}. Overlay networks such as Tailscale, which runs over
WireGuard~\cite{donenfeld2017wireguard}, solve the same problem with a
coordination service and relays; PipedPeer embeds the coordination in its own
daemons and never relays job data.

\paragraph{Scheduling on heterogeneous machines} Dividing divisible work among
machines of different speeds to minimise makespan is classical scheduling on
uniform machines~\cite{graham1969bounds,hochbaum1988uniform}. MapReduce's
backup tasks~\cite{dean2004mapreduce} and LATE~\cite{zaharia2008late} launch
speculative copies of stragglers; PipedPeer applies the same idea to the parts
of one parallel map. Edge-computing work studies the offloading decision our
cost model makes: weighing transmission against local
computation~\cite{deng2023}, partial over all-or-nothing
offloading~\cite{yong2023}, learned policies~\cite{he2024}, and
topology-aware placement~\cite{zhao2023}, which outperforms the scalar scoring
PipedPeer uses.

\paragraph{Data-parallel training over weak links} PyTorch
DDP~\cite{li2020pytorchddp} and Horovod~\cite{sergeev2018horovod} use ring
all-reduce among ranks that reach one another directly. Communication is
reduced by 8-bit~\cite{dettmers2016eightbit} and 1-bit~\cite{seide2014onebit}
gradient compression with error feedback, by low-rank
compression~\cite{vogels2019powersgd}, and by synchronising less often: local
SGD~\cite{stich2019localsgd,lin2020localsgd}, federated
averaging~\cite{mcmahan2017fedavg} and DiLoCo~\cite{douillard2023diloco}.
Collaborative training over the internet has been demonstrated by
Learning@home~\cite{ryabinin2020learningathome}, DeDLOC~\cite{diskin2021dedloc}
and Petals~\cite{borzunov2023petals}. PipedPeer's training path is smaller in
ambition than these: two to four ranks of an unmodified script, synchronised
through the daemons, with the synchronisation interval chosen from accuracy
measurements (Section~\ref{sec:eval-train}).

\paragraph{Failure transparency} Spenger~\cite{kth2025thesis} defines failure
transparency, under which application code never observes a failure. Scripts
run by PipedPeer contain no failure handling, so every distributed path must
fall back to a local one (Section~\ref{sec:ft}).

Table~\ref{tab:cmp} summarises the comparison. Its entries come from published
designs, not from benchmarks.

\begin{table*}[t]
\centering
\caption{PipedPeer against existing systems (qualitative). The lower block is
where PipedPeer is partial or the weaker system.}
\label{tab:cmp}
\footnotesize
\begin{tabularx}{\textwidth}{X C{1.9cm} C{1.1cm} C{1.1cm} C{1.1cm} C{1.4cm} C{1.0cm}}
\toprule
\textbf{Property} & \textbf{PipedPeer} & \textbf{Spark} & \textbf{Ray} & \textbf{Dask} & \textbf{Lithops} & \textbf{ssh} \\
\midrule
Runs unmodified Python     & yes & no  & partial & partial & partial & yes \\
No cluster or back end to configure & yes & no & no & no & no & partial \\
No head or master node     & yes & no  & no      & no      & n/a     & yes \\
Reproduces the environment & yes & partial & partial & partial & partial & no \\
Joins NATed machines, no relay$^{\dagger}$ & yes & no & no & no & n/a & no \\
\midrule
Survives losing a node$^{*}$ & partial & yes & yes & yes & yes & no \\
Authenticated, encrypted   & partial & yes & partial & partial & yes & yes \\
Sandboxes each job         & partial & no  & no      & no      & yes & no \\
Runtimes beyond Python     & no  & yes & partial & no      & no      & yes \\
Proven in production       & no  & yes & yes     & yes     & yes     & yes \\
\bottomrule
\end{tabularx}
\par\smallskip\raggedright\footnotesize
$^{\dagger}$When one introducer is reachable; otherwise a public address or an
overlay network is needed.
$^{*}$Any worker; not the submitter of a parallel script nor the machine of a
training ring's rank~0.
\end{table*}

% ============================================================
\section{Design overview}
\label{sec:design}

\paragraph{Goals}
\textbf{(G1)} The script runs unchanged.
\textbf{(G2)} Joining needs one setup command and, across networks, one invite
code.
\textbf{(G3)} No machine is special: there is no head node, scheduler or
registry, and every machine both submits and works.
\textbf{(G4)} Distribution never changes the result and is declined when the
model predicts it would be slower.
\textbf{(G5)} No job data passes through a third machine.

\paragraph{Architecture} Every machine runs a single daemon; the command-line
client is short-lived (Fig.~\ref{fig:arch}). The coordinator is a library
linked into the client, so placement is decided where the job is submitted.
Every daemon also serves the introducer for its own cluster, and a training
ring's lead is whichever node runs rank~0. On one network participants find
each other by broadcast; across networks a participant runs
\texttt{pipedpeer connect} with a one-line code printed by another member's
\texttt{pipedpeer invite}, which carries the cluster secret and up to six
introducer addresses with a checksum.

\begin{figure*}[t]
\centering
\includegraphics[width=0.75\textwidth]{\figpath{fig_3.1_system_design.png}}
\caption{System design. Submitter and worker are the same binary. Peers on
other networks are reached over direct links that the introducer helps set up
but never carries. Phones take parallel-map and torch chunks only.}
\label{fig:arch}
\end{figure*}

\paragraph{Life of a job} (1) The client scans the script's imports and builds
or reuses its closure. (2) It places the job and leases a slot on the chosen
worker. (3) It sends only the store paths the worker lacks, then the workspace.
(4) The worker runs the script in a sandbox with the interception package on
the interpreter's path. (5) Parallel calls inside the script may be divided
among peers that can run the same closure. (6) Output is streamed back and
changed files are synchronised to the submitter.

% ============================================================
\section{Mechanisms}
\label{sec:mech}

\subsection{Placement}
\label{sec:placement}

Candidates come from known peers, LAN discovery and internet-linked peers.
Hard filters run in order --- architecture, GPU if required, then free memory,
cores and VRAM --- and the remaining candidates are scored:
\begin{equation}
\label{eq:score}
\begin{split}
S = \Bigl(1 &- \frac{u_{\mathrm{cpu}} + u_{\mathrm{mem}}}{200} - 0.05\,n_{\mathrm{jobs}} \\
    &+ 0.10(\hat{c}-0.5) + 0.06(\hat{f}-0.5) \\
    &+ 0.10(\hat{m}-0.5) + 0.10(\hat{r}-0.5)\Bigr)\cdot h ,
\end{split}
\end{equation}
with utilisations $u$ in percent, $n_{\text{jobs}}$ running jobs, capacity
terms $\hat{c}=\min(\text{cores}/32,1)$ and $\hat{f}=\min(\text{MHz}/4000,1)$,
headroom terms $\hat{m}=\min(\text{free GiB}/32,1)$ and $\hat{r}$ (free fraction
of cores), $0.5$ for a missing metric, and $h$ a health factor that halves each time the
node fails a job or stops answering (floor 0.1) and regains 0.1 per healthy
poll, up to 1.0, so a flaky node loses to an equal healthy one.
% source: commit 23137d0 (nodestore.MarkFailed, HealthFloor, HealthRecovery). $S$ is clamped at zero; a GPU request adds up to 0.8 from
VRAM, device count and compute capability. The weights are hand-set: the
utilisation terms can subtract up to 1.0 and the four normalised terms move $S$
by at most 0.18 either way, so current load dominates hardware. We did not run
a sensitivity analysis of these weights (Section~\ref{sec:threats}).

Free memory is what a node can \emph{grant}: where a user \texttt{systemd} is
available, each daemon runs under a cgroup ceiling of half the machine, and
daemons co-located on one machine subtract one another's reservations. Work is
reserved by a three-phase lease (accept, commit, complete); the concurrency
and memory checks and the insertion happen under one lock. The submitter renews
every 30\,s, also while the environment builds and uploads, against a 5-minute
expiry. A submitter that dies stops renewing and its slot is reclaimed; a
worker that dies has its job rescheduled with back-off.

\subsection{Environment reproduction}
\label{sec:env}

The client resolves the script's imports to packages and builds a Nix closure
from a flake pinned to an \emph{exact} package-set revision, so the store path
is a stable name for an environment across machines and time. Imports absent
from the mapping table are looked up by name in the pinned set, the scan
follows the script's own modules, a \texttt{requirements.txt} is honoured, and
an import that is still unresolved is reported before the run.

Before uploading, the client asks whether the target can already run the
closure and otherwise sends only the store paths it lacks. The worker then
pushes the closure to peers that cannot yet run it, since only peers that can
are offered distributed work. Between daemons the push is a binary cache
produced and imported by Nix itself and signed with the sending node's key;
the receiver checks that key against the signer's published identity before
trusting it for that import. The upload returns once the target holds the
closure; pushes to slower peers continue in the background and do not hold up
the run.\footnote{An earlier build made the upload wait on the synchronous
push to every peer. In an emulated three-node lab this blocked a cold run for
about 600\,s and, because a peer still importing was asked again, a warm run
too (579\,s). The behaviour described here replaced it
(Section~\ref{sec:eval-cache}).}
% source: docs/benchmarks/scaling.md "cold"; fix 1f10629; after-fix numbers commit a56e4de.
Closures go only to peers of the same architecture; a peer of another
architecture receives the job's package list and builds its own environment
once, caching it by flake hash.

\subsection{Joining across networks}
\label{sec:nat}

\begin{figure}[t]
\centering
\includegraphics[width=\columnwidth]{\figpath{fig_4.1_nat_traversal.png}}
\caption{Joining across networks. The introducer exchanges candidate addresses
and forwards one punch request; the two nodes then race their candidate pairs
over one UDP socket.}
\label{fig:nat}
\end{figure}

An introducer keeps an address book keyed by a hash of the cluster secret and
forwards punch requests between two nodes; it carries no job data. Every daemon
serves one for its own cluster and registers with every introducer it knows,
so no single machine's address is a point of failure. Each node offers its LAN
addresses, a port mapped on its router (published only when the router's
external address matches the one the introducer observed, which excludes a
router behind a second NAT) and its reflexive address (Fig.~\ref{fig:nat}).
Probes, introducer messages and QUIC share one UDP socket, because a router's
mapping belongs to the socket that opened it; both sides probe every candidate
pair every 250\,ms for 15\,s rather than trying candidates in turn.

When one side's NAT allocates a fresh external port per destination, there is
no address to publish. That side opens 256 sockets towards the peer while the
other sends 600 probes to random ports among the 64{,}511 above 1024, and the
socket that is hit is kept~\cite{anderson2020nat}. If the NAT draws ports
uniformly, a round succeeds with probability
\begin{equation}
\label{eq:birthday}
P = 1-\left(1-\frac{256}{64511}\right)^{600} \approx 0.91 .
\end{equation}
This is a computed figure under that assumption; we did not measure a
per-round success rate. A container lab with two cone and two port-randomising
translators exercises the fallback (cone--randomising pairs are expected to
link through it, and a pair of two randomising NATs to be reported unreachable),
and the four defects it exposed --- among them a socket found by a collision
that stopped listening once adopted --- are fixed.
% source: commits 2e3e344 (lab/nat), 8944270 and docs/system-report.md "What the NAT lab found".

The winning path carries QUIC with a hello signed by each node's
Ed25519~\cite{bernstein2012ed25519} key and checked against the node it meant
to reach. The hello signs the session's exported keying
material~\cite{rfc5705} and proves the cluster secret by an
HMAC~\cite{rfc2104} over it, so it cannot be replayed on another session. The
link appears locally as a loopback forwarder recorded as one path to that node,
so placement, transfer and gradient exchange run over it unchanged. A node is
identified by its key's fingerprint, never by its address: the node store keeps
every path to a node and switches only to a path that answered and proved the
key, and QUIC connections migrate when an address changes. A peer with no path
is reported with its reason and routed around, never relayed. The daemon API
requires the cluster secret on every route and serves pinned
TLS~\cite{rfc8446} and HTTP on one port.

\subsection{Transparent interception}
\label{sec:intercept}

\begin{figure}[t]
\centering
\includegraphics[width=\columnwidth]{\figpath{fig_4.8_shim_injection.png}}
\caption{Interception. The package is placed on the interpreter's path so that
it is imported before the script; libraries are patched on first import.}
\label{fig:shim}
\end{figure}

A Python package embedded in the binary is placed on the interpreter's path so
that it is imported before the user's first line, and libraries are patched on
their first import rather than at startup (Fig.~\ref{fig:shim}). Substituted
primitives are \texttt{multiprocessing} and \texttt{concurrent.futures} maps,
\texttt{joblib}, pandas grouping, joining and chunked reading, NumPy matrix
products and decompositions, and PyTorch data-parallel training. Every
substituted call falls back to the original on any exception.

A parallel-map kernel travels as \emph{source} when it is a module-level
function, including one defined in \texttt{\_\_main\_\_}, or \emph{by value}
through cloudpickle~\cite{cloudpickle} when it is a lambda or closure; a kernel
that cannot be serialised runs on local threads. pandas grouping and joining
hash rows by key, so each key is complete on one node and any aggregation is
exact without a combiner; the row order of a join is not preserved.

\subsection{Dividing a parallel map}
\label{sec:split}

The node that holds a map splits it once among its peers, and each peer runs
its part on its own cores; a part is never split across machines again. Each
device $i$ (a node's CPU pool or one GPU) is modelled by a setup time $o_i$ and
an effective rate
\begin{equation}
\label{eq:rate}
r_i = \Bigl(\frac{1}{\lambda_i} + \frac{\beta}{B_i}\Bigr)^{-1},
\end{equation}
where $\lambda_i$ is its measured throughput in items per second, $\beta$ the
bytes per item and $B_i$ the measured bandwidth to it ($r_i=\lambda_i$ for the
local device). A share of $w_i$ items finishes at $o_i + w_i/r_i$. For $W$
items, the makespan $T$ is the root of
\begin{equation}
\label{eq:makespan}
\sum_i \max(0,\,T-o_i)\,r_i = W ,
\end{equation}
which is piecewise linear and increasing in $T$ and is solved by walking its
breakpoints. A device with $o_i \ge T$ cannot start before the others finish
and is dropped; dropping it raises $T$, so the rule is applied until the set
settles (Algorithm~\ref{alg:split}). Rates and setup times are fitted from the
parts each device has returned. Because the local device is always a
candidate, the \emph{predicted} makespan is never above the local one; the
measured time can be (Section~\ref{sec:eval-dist}).

\begin{algorithm}[t]
\caption{Dividing $W$ items among devices}
\label{alg:split}
\begin{algorithmic}[1]
\State $D \gets$ devices with a measured rate and enough memory for the working set
\Repeat
  \State $T \gets$ root of $\sum_{i\in D}\max(0,T-o_i)\,r_i = W$
  \State $X \gets \{\, i \in D : o_i \ge T \,\}$; \quad $D \gets D \setminus X$
\Until{$X = \emptyset$}
\State $w_i \gets \lfloor (T-o_i)\,r_i \rceil$ for $i\in D$; return shares below a minimum to the others
\end{algorithmic}
\end{algorithm}

During execution a part that runs past its prediction counts as at least that
long, and a part that is waiting or overdue is copied to a free device when the
expected saving exceeds the copy's own setup (or 0.2\,s); the first copy to
finish is kept, in the manner of backup tasks~\cite{dean2004mapreduce,zaharia2008late}.

\subsection{Cost model for numerical work}
\label{sec:cost}

NumPy and pandas operations are gated. With $\phi$ operations and $\rho$ link
crossings per byte, bandwidth $B$, local and per-node remote throughput $R$ and
$R'$, and $K$ nodes holding the environment (the submitter included), divisible
work moves only when its time per byte is lower on the cluster:
\begin{equation}
\label{eq:cost}
\frac{\phi}{R} \;>\; \frac{\rho}{B} \;+\; 1.3\,\frac{\phi}{R'K} .
\end{equation}
Payload size cancels, so size enters as separate rules: at least two nodes and
32\,MiB, and pandas work below 8\,FLOP/B (grouping $\phi=4$, joining $\phi=2$,
$R=1$\,GFLOP/s) stays local up to 512\,MiB. A matrix product
($\phi=\max(8,n/8)$, $\rho=5$, $R=200$\,GFLOP/s) that does not fit in memory
moves regardless, or raises an error if it fits nowhere. One GPU peer sets
$R'=6R$ for all $K$ nodes. The factor 1.3, the GPU factor and the throughput
constants are coarse, hand-set values that we did not calibrate across
hardware. The test is first evaluated at 10\,Gbit/s, and bandwidth is measured
only if offloading could then win (Fig.~\ref{fig:cost}). Indivisible work
(singular value decomposition) moves when its transfer costs under half the
local time. On the machines tested, local BLAS outruns a CPU cluster's
transport, so dense products stay local.

\begin{figure}[t]
\centering
\includegraphics[width=\columnwidth]{\figpath{fig_4.11_cost_model.png}}
\caption{Decision regions of Eq.~\eqref{eq:cost} for pandas work ($\rho=1$,
$R=R'=1$\,GFLOP/s) with three nodes on a 500\,MB/s link. Above 512\,MiB,
work above about 3.5\,FLOP/B is distributed: a 900\,MiB grouping moves, a join
of the same size stays local.}
\label{fig:cost}
\end{figure}

\subsection{Data-parallel training}
\label{sec:train}

\begin{figure}[t]
\centering
\includegraphics[width=\columnwidth]{\figpath{fig_4.10_ddp.png}}
\caption{Training through the daemons. Ranks post gradients (or weights) to
daemons rather than to one another and receive one sample-weighted average.}
\label{fig:ddp}
\end{figure}

Ranks post gradients to daemons rather than to one another
(Fig.~\ref{fig:ddp}), which avoids the rank-to-rank reachability that ring
all-reduce assumes. Each rank posts to the lead node's daemon; in exact mode on
three or more directly reachable machines, it posts one shard to each rank's
daemon. Every rank receives one sample-weighted average
$\sum_i n_i g_i / \sum_i n_i$, sent as 8-bit values with per-block scales and
error feedback~\cite{dettmers2016eightbit,seide2014onebit}. The global batch is
partitioned across ranks in proportion to their measured speed, preserving the
effective learning rate.

Unless the user fixed the ring size with \texttt{-{}-ddp~$N$}, after a warm-up
the lead compares the split's step time with the fastest rank alone and keeps
the ring only if it is faster. Admission takes the $k$ fastest candidates
maximising $k\,r_k$, the equal-batch throughput; this is conservative, since
placement cannot price synchronisation. Samples go to ranks by a weighted
rendezvous hash~\cite{thaler1998rendezvous}: when a rank dies the others take
its share at the next epoch boundary with over 95\% of their samples in place,
and a returning rank joins from the lead's snapshot, taking about $1/n$ of them.

A fast mode synchronises every $H$ steps: an exact warm-up (a tenth of the
run, 10--500 steps), then local steps whose weights are averaged by samples,
with optimiser moments averaged every tenth round; $H$ shrinks to 1 over the
run's last tenth. Fast mode never changes epochs, batch or learning rate, and
must keep validation loss at most 1\% above one device's (mean over seeds, on
every test model). DiLoCo's outer Nesterov step (learning rate 0.7, momentum
0.9) multiplies each round's averaged change by $0.7/(1-0.9)=7$ in the steady
state; DiLoCo synchronises every few hundred steps, and at our intervals of
2--64 steps the amplified update failed the 1\% bar, so PipedPeer uses plain
averaging (outer learning rate 1, no momentum) and chooses $H$ by ring size
(Section~\ref{sec:eval-train}). Automatic mode never picks fast mode beyond
four ranks; an explicit fast ring above four runs at $H=4$ with a warning.
Under automatic synchronisation the lead switches a ring to fast mode once,
when the predicted fast step with a 10\% margin is shorter than the exact
split's and, for a ring not fixed by \texttt{-{}-ddp~$N$}, also shorter than
the fastest rank's alone.

\subsection{Degradation and operation}
\label{sec:ft}

Using the cluster must not change the result. Parallel-map parts fill a result
table in which the first result for each part is kept, idle local cores re-run
any part not yet returned, and the call returns once every item is covered, so
an unresponsive cluster degrades to a local pool. A failed peer is skipped, a
dead worker process is respawned, a failed job is rescheduled with back-off,
and a training ring continues on the ranks that are left. Two machines remain
single points of failure: the submitter of a parallel script, and the machine
of a training ring's rank~0.

A run behaves like a local one. An interrupt sends one cancel frame, which
stops the job's process group and its parts on every peer. A dropped
connection never kills the job: the daemon logs every output frame, and the
submitter can reattach and replay it. Each run ends with a line splitting its
time between the script and PipedPeer, and with where its items ran.

\subsection{Phones}
\label{sec:phone}

An Android phone joins from Termux~\cite{termux} without root. Setup installs a
Debian root through \texttt{proot-distro}, which uses
\texttt{proot}~\cite{proot}, with single-user Nix inside it, and falls back to
Termux's own Python when that fails. The phone advertises a
\emph{chunks-only} role: it takes parallel-map and torch chunks but not whole
jobs or training ranks, because the job sandbox would need a second
\texttt{proot} nested in the first, which the Termux build we tested cannot
start. Since closures go only to peers of the same architecture, a phone
receives the job's package list, builds an aarch64 environment (mostly fetched
from the public binary cache), and accepts the job's chunks once the build is
done. Android offers apps neither user namespaces nor cgroups, so the phone's
per-job isolation is a \texttt{proot} view of the job directory and the store
with an address-space limit and a lower CPU priority, which is weaker than the
\texttt{crun} sandbox used elsewhere.

\subsection{Security model}
\label{sec:security}

The trust boundary is the cluster: anyone holding the shared secret is a
member. Links are authenticated by node keys and the secret, closures are
signed by the sending node, and on Linux each job runs in an unprivileged
\texttt{crun}~\cite{crun} container with a memory ceiling, without network
isolation. Three gaps remain: TLS on the daemon API is trust on first use; the
keys trusted to sign closures are not yet bound to an authenticated identity;
and authorisation is one shared secret, so any member can run code on any
other. PipedPeer is therefore suited to machines whose owners trust one
another, not to open volunteer computing.

% ============================================================
\section{Implementation}
\label{sec:impl}

PipedPeer is one Go~1.25 binary, client and daemon, of 46 Go packages, with
the interception package embedded as Python sources. It links peers with
quic-go, builds environments with Nix flakes, and runs each job in a
\texttt{crun} sandbox. The Go test suite, run with NumPy, pandas,
scikit-learn, joblib and torch installed, passed 1{,}012 of its 1{,}017
top-level tests at commit 6bc0ef0. Four were skipped (two need a container
runtime, one a Termux environment, one a system-wide Nix), and one failed: a
test requiring a cancelled job to stop within one second took 1.56\,s; it also
failed one of five isolated reruns, so it is timing-sensitive.
% source: go test -count=1 -json ./... at 6bc0ef0, 2026-11-02, torch 2.15.0+cpu.
A separate
check fails the build if any of the tests that exercise distribution was
skipped, so a missing library cannot turn them into silent passes. A mutation
audit disables each of 95 features in turn --- among them the sample-weighted
gradient average, the refusal to publish a router mapping behind a second NAT,
and the check that a signing key belongs to the node offering it --- and
requires a named test to fail; all 95 were caught~\cite{jia2011mutation}.
% source: docs/benchmarks/raw/T3_audit_tests.txt (27f94cd, run on 8d63c30).
The source code and every measurement record are public~\cite{pipedpeerrepo}.

% ============================================================
\section{Evaluation}
\label{sec:eval}

We ask five questions.
\textbf{Q1:} What does interception cost, including when it declines?
\textbf{Q2:} How much do the environment cache and the closure diff save?
\textbf{Q3:} Does work cross machines with correct results, how does the
division behave on unequal machines, and how does cost grow?
\textbf{Q4:} Does split training keep accuracy, and when is it faster?
\textbf{Q5:} What happens between physical machines on different networks,
including a phone?

\subsection{Setup and method}
\label{sec:setup}

We use four settings and label every result with its setting.

\emph{Plain lab.} Three worker daemons in containers on one host, sharing its
processor, memory and loopback interface. It measures correctness, overhead
and caching; it cannot measure speedup, because every container competes for
the same processor.

\emph{Quota lab (emulated).} Each container gets a fixed CPU and memory quota
and links shaped with \texttt{tc netem} to a LAN (0.25\,ms one way,
1000\,Mbit/s), Wi-Fi (4$\pm$2\,ms, 80\,Mbit/s) or 5G (15$\pm$5\,ms,
100\,Mbit/s). Speedup is one container against several of the same size.
Results from this setting are \emph{emulated}: they show how cost grows with
one variable at a time, not what physical machines achieve. Some sweeps ran
while other work loaded the host (one-minute load 0.4--4 on 15 cores); their
growth fits use CPU time, which load barely moves, and their wall times are
marked.

\emph{Two-site testbed.} Two containers on separate private networks, both
numbered 192.168.0.0/24, each behind a router container that translates
addresses and drops unsolicited inbound packets, with the introducer on a third
network standing for the internet. A script runs every command of the
project's manual test plan across it and checks each exit code and, where work
should cross, that both sites took a share.

\emph{Physical machines.} Two sets. \emph{Set A}: two laptops in different
homes --- an AMD Ryzen~7 5800H with an RTX~3060 and about 14\,GiB, and an Intel
Core i5-13420H with an RTX~3050 and 15.3\,GiB --- joined over a direct Tailscale
path with a 6--8\,ms round trip, plus a Samsung Galaxy S20~FE phone (Android~13)
under Termux. Neither home offered a reachable introducer (one router refuses
port mapping, the other laptop sits behind two NATs), so Set~A measures
distribution, not NAT traversal. \emph{Set B}: a 20-core 13th-generation Intel
Core i9 machine and the Ryzen laptop over Tailscale, training over a home link
of about 54\,Mbit/s, and a public cloud host serving as introducer.

The labs ran on the two laptops of Set~A (Fedora~43 and~44). Lab wall times are
medians of three to five runs; training rings are one seed unless stated;
interpreter startup is the best of five; sandbox start is the mean of 20 runs
after five warm-up runs. Physical-machine figures are single runs or ranges
over two or three runs, as stated. Accuracy is the mean over seeds against one
device trained on the same data and seeds. MiB is binary, MB decimal.

\paragraph{Test models} Three models (Table~\ref{tab:models}) keep their batch,
learning rate and epochs in every arm.

\begin{table*}[t]
\centering
\caption{Test models for training accuracy}
\label{tab:models}
\footnotesize
\begin{tabularx}{\textwidth}{L{2.0cm} X C{2.1cm} C{1.4cm} C{1.0cm} C{1.6cm}}
\toprule
\textbf{Model} & \textbf{Network and optimiser} & \textbf{Parameters} & \textbf{Samples} & \textbf{Batch} & \textbf{Epochs (steps)} \\
\midrule
MLP-G & MLP 64--512--512--512--1; SGD 0.05, momentum 0.9 & 559{,}105 & 200{,}000 & 4096 & 3 (147) \\
MLP-C & Same MLP; SGD 0.1, momentum 0.9 & 559{,}105 & 60{,}000 & 2048 & 6 (180) \\
Transformer & 2 layers, width 64, 4 heads, 32 tokens; AdamW & 70{,}211 & 40{,}000 & 256 & 2 (314) \\
\bottomrule
\end{tabularx}
\end{table*}

\subsection{Q1: Interception cost}
\label{sec:eval-shim}

\begin{figure*}[t]
\centering
\includegraphics[width=0.8\textwidth]{\figpath{fig_5.2_shim_overhead.png}}
\caption{Cost of interception. Left: a declined 16-item map, median and range
of five runs (the enabled side's cold first run is excluded from the range).
Right: interpreter startup with libraries patched at startup and on first
import, best of five.}
\label{fig:shimcost}
\end{figure*}

On the Core~i5 laptop, startup of a job that imports nothing heavy took
0.117\,s with interception against 0.094\,s without, a ratio of $1.24\times$
(best of five; Fig.~\ref{fig:shimcost}). An earlier design that patched NumPy,
pandas, torch and joblib at startup measured $20.0\times$, on a different
machine. A continuous-integration check fails any change that raises the ratio
above $3\times$.
% source: docs/benchmarks/raw/O2_shim_startup.txt; B3_shim_overhead.txt.

A 16-item map whose tail is 0.01\,s of work, below the 0.5\,s dispatch floor,
is declined by the cost model. Its median over five runs was 0.796\,s with
interception against 0.678\,s without, every run's receipt showing no item
dispatched: declining costs about 0.12\,s per job. A forced run of the same
workload did dispatch, so the decline is the model's choice and not an
unavailable cluster.
% source: docs/benchmarks/RESULTS.md R3 [Oct-12], raw/N7_*.

The \texttt{crun} sandbox costs 15--21\,ms more per job than Bubblewrap across
three tasks (means of 20 runs), small beside a warm job of under a second;
its device model is what enables GPU passthrough.
% source: docs/benchmarks/raw/B2_sandbox_bench.md [Aug].

\subsection{Q2: Environment cache and closure diff}
\label{sec:eval-cache}

\begin{figure}[t]
\centering
\includegraphics[width=\columnwidth]{\figpath{fig_5.1_closure_cache.png}}
\caption{Five submissions of the same environment in the plain lab.}
\label{fig:cache}
\end{figure}

The same trivial job was submitted five times to the plain lab
(Fig.~\ref{fig:cache}). The first submission, which also built the
environment, took 23.4\,s. The second landed on a worker that did not yet hold
the closure and took 8.1\,s. The next three found it in place, sent only the
workspace, and took a median of 0.83\,s.
% source: docs/benchmarks/RESULTS.md R2, raw/O1_bench_py.txt.

Against a container peer on the same host, adding one package to an
environment the peer held sent 7 of 43 store paths, 28.6\,MB instead of
206.0\,MB (86\% fewer bytes); a peer fetched 0.45\,MB of a 466\,MB closure
that differed from one it held by one small pure-Python package. In the LAN
lab, a dry-run build of an unseen script's environment took 4\,s on a container
that had built the common packages ahead of time, against 54\,s on one that had
not.
% source: commit 193e1bf (raw/C_commit_measurements.md); rehearsal.md row "Teacher-style script's environment".

\paragraph{A defect found by measurement} In an emulated three-node lab, the
first version of closure broadcast blocked a cold run for 603\,s and a warm run
that followed it for 579\,s: the upload's reply waited on a synchronous push
to every peer, and a peer still importing the last closure was asked again
until the 10-minute client timeout. After the push was moved to the background
(Section~\ref{sec:env}), the same lab took 137.5\,s cold and 62.4\,s warm. With
one of two nodes held to 800\,kbit/s, the first run waited the upload's 60\,s
limit for the slow node and then started; once the background send finished, a
run took 2.1\,s with the slow node taking a share.
% source: docs/benchmarks/scaling.md "cold"; commit a56e4de body; report results.tex sec:res-cache.

\subsection{Q3: Distribution, division and growth}
\label{sec:eval-dist}

\paragraph{Correctness} With distribution forced in the plain lab, the node
running a 48-item parallel map sent 18 items to three other daemons in each of
five runs, with results identical to a local run. A NumPy workload moved its
34\,MB singular value decomposition to one worker. Unforced, a 48-item pool
benchmark took 0.17\,s in plain Python, 0.17\,s with a hand-tuned pool and
0.24\,s under PipedPeer with 12 items on the cluster: on one shared host,
distribution can only add overhead, and here it added 0.07\,s.
% source: RESULTS.md R4, raw/N3_*, raw/N5_*, raw/O3_bench_pool_lab.txt.

\paragraph{Division on unequal machines} On a quota lab of 8, 2 and 1 CPUs,
dividing a map by headcount took 3.6\,s, $1.6\times$ slower than not
distributing (which puts the local run near 2.3\,s). Dividing by measured
throughput declined the 1-CPU worker and took 2.9\,s: better than the
headcount split, but still about $1.3\times$ slower than local. The model
predicted a gain that the measurement did not show; at this size the map is
better kept on one machine. With one peer held to a tenth of a CPU in a cluster
of otherwise equal peers, the slowdown against an equal cluster fell from
$1.6\times$ on the first run to $1.08\times$ by the fourth, as the rate
estimates converged.
% source: commit 8c9284c (raw/C_commit_measurements.md); report results.tex "Sharing by measured speed".

\paragraph{How cost grows (emulated)} A sweep varied one quantity at a time
around 64 items of about 75\,ms and 256\,KB each on nodes of 2 CPUs and
768\,MiB: payload 1\,KB--64\,MB, 8--2048 items, per-item cost
$10^3$--$3\times10^6$ loop iterations, 2--4 nodes, and a speed spread of up to
$10\times$ (Table~\ref{tab:scaling}). No step's CPU time grew faster than its
expected order (log--log slope above the expected exponent by more than 0.15).
Below about $10^5$ iterations per item the map stayed local, as designed. One
peer at a tenth of the speed raised wall time $1.9\times$ where lost capacity
alone predicts about $1.4\times$, a straggler tail. Five percent packet loss
moved work back to the submitter (peers 3.0 to 1.7 CPU-s) and raised wall time
from 1.6 to 2.9\,s; 1--2\% cost little. A 256\,MB map failed on a 768\,MiB
submitter for lack of memory.
% source: docs/benchmarks/scaling.md (binary fc83359, 3 reps per cell, medians).

\begin{table}[t]
\centering
\caption{Emulated node-count sweep: a 2.5 CPU-second map on equal 2-CPU nodes,
median of three runs on a host also carrying other work (one-minute load in
the last column), so wall times are inflated unevenly. Not a physical-machine
speedup.}
\label{tab:scaling}
\footnotesize
\begin{tabular}{ccccc}
\toprule
\textbf{Nodes} & \textbf{Wall (s)} & \textbf{vs.\ one node} & \textbf{Total CPU (s)} & \textbf{Host load} \\
\midrule
1 & 2.37 & $1.0\times$ & 2.54 & 2.53 \\
2 & 1.15 & $2.1\times$ & 3.50 & 0.96 \\
3 & 0.95 & $2.5\times$ & 3.55 & 0.44 \\
4 & 0.81 & $2.9\times$ & 3.34 & 0.66 \\
\bottomrule
\end{tabular}
\end{table}

The step from one node to two adds about 1\,CPU-s (2.5 to 3.5), the fixed cost
of distributing; from two to four nodes total CPU stays flat.

\paragraph{End to end} Table~\ref{tab:e2e} gives whole-system runs in a LAN lab
of four containers on the Ryzen laptop (and a two-container variant sharing its
GPU). The containers share one machine, so these are costs, not speedups. Cold
times are dominated by building and fetching environments, the GPU steps by
downloading CUDA PyTorch.

\begin{table}[t]
\centering
\caption{End-to-end steps in the LAN lab, first run and cached (one host)}
\label{tab:e2e}
\footnotesize
\begin{tabularx}{\columnwidth}{X C{1.7cm} C{1.7cm}}
\toprule
\textbf{Step} & \textbf{First run} & \textbf{Cached} \\
\midrule
Parallel map & 43.5\,s & 17.0\,s \\
Map over a lambda & 20.5\,s & 4.8\,s \\
GPU job submitted from a node without a GPU & ---$^{a}$ & 7.6\,s \\
Two-rank ring on one shared GPU & 630.3\,s & 9.5\,s \\
Script combining every intercepted kind of work & 664.6\,s & 26.7\,s \\
Five idiomatic user scripts$^{b}$ & 35--126\,s each & 20--70\,s each \\
Eight jobs at once (6.0\,GB of memory in total) & \multicolumn{2}{c}{412.0\,s, all eight completed} \\
\bottomrule
\end{tabularx}
\par\smallskip\raggedright\footnotesize
$^{a}$The recorded first run measured a closure broadcast to nodes without a
GPU, which a later change removed, so it is not reported.
$^{b}$First run after the common packages were built ahead of time.
\end{table}
% source: docs/benchmarks/rehearsal.md "End-to-end steps".

\paragraph{Failure drills} Table~\ref{tab:drills} lists drills that remove
machines, links and terminals during a run. They exercise rescheduling,
re-sharding at an epoch boundary, the cancel frame and the logged output
stream. A fault-injection run in the plain lab first tested nothing, because
no worker held the closure and the map stayed local; once every upload seeded
the peers, the same harness passed (second row). The test suite adds that a parallel map whose every remote part fails still
returns the correct result.
% source: RESULTS.md R5 (failed in October; passed per commit a9c715d); rehearsal.md "Failure drills" and "Two-site testbed".

\begin{table*}[t]
\centering
\caption{Failure drills (LAN lab unless stated)}
\label{tab:drills}
\small
\begin{tabularx}{\textwidth}{X L{6.4cm}}
\toprule
\textbf{Drill} & \textbf{Result} \\
\midrule
A node's link cut during a 240-item map, two-site testbed & its 53 items re-run on the submitting node once the part overran its predicted 18\,s (detected at 81\,s); checksum unchanged; map took 83.1\,s \\
A worker killed with map chunks in flight (plain lab, 80 items)$^{d}$ & 33 items had run on peers; no chunk lost; result correct \\
A worker killed during a remote job & node reported gone 16\,s after the kill \\
A training rank killed during a run & remaining rank finished all 80 epochs, 90.5\,s \\
A training rank's machine back during a run & rejoined at the next epoch, 4\,s after its daemon started \\
Interrupt during a parallel map & peers stopped within 2\,s \\
Terminal closed during a run & stopped everywhere, as an interrupt \\
Detach, then attach during and after the run & followed live; replayed in 1.0\,s once finished \\
Submitter's network removed for 200\,s during a remote job & job ran once; result back about 15\,s after reconnecting \\
Introducer's address changed, no command run$^{c}$ & links held or re-formed with no command; reported 12.1--12.6\,s \\
\bottomrule
\end{tabularx}
\par\smallskip\raggedright\footnotesize
$^{c}$The harness counts a recovery no earlier than 12\,s after the move when
the path survives outright, so these readings sit at its floor and do not
measure a recovery time.
$^{d}$Recorded in commit a9c715d (3 November 2026); its raw output is not stored
in the repository.
\end{table*}
% source: lab-net.sh recover() (12 s floor); commit c785275 body (12.1-12.6 s).

\subsection{Q4: Training}
\label{sec:eval-train}

\paragraph{When splitting pays} In the quota lab, a small model with heavy
steps on a LAN trained in 9.1\,s split across two 2-CPU containers against
11.8\,s on one, while a 9\,MB model with light steps on Wi-Fi took 80.7\,s split
against 8.0\,s on one container. A GPU widens the gap: on the Ryzen laptop one
training job took 1.5\,s on its RTX~3060 against 55.0\,s on its CPU (one run).
% source: docs/benchmarks/ddp-switch.md (quiet rerun, a56e4de); commit 2da0857.

\paragraph{Fast synchronisation} At the default two-rank interval of 32 steps,
in-process runs kept validation loss within 1\% of one device on every model:
MLP-G $+0.06\%$, MLP-C $+0.64\%$ and the transformer $+0.73\%$ (seeds 0--2),
and MLP-C $+0.35\%$ and the transformer $-1.02\%$ over seeds 0--5 (seeds 3--5
were run for those two only). Single seeds spread from $-3.9\%$ to $+2.2\%$
across the runs reported here, so the bar holds for the mean, not for every
seed. Table~\ref{tab:fastsync} gives one lab run (one seed) at this interval:
each rank sent 7.6--9.2$\times$ fewer bytes than under exact synchronisation,
and the ring finished 1.2--1.4$\times$ faster than exact on the LAN profile and
1.6--3.1$\times$ on Wi-Fi (Fig.~\ref{fig:sync}). Against one container, the
fast ring was faster on every model and profile here (for example MLP-G
15.2\,s on LAN against 27.9\,s), because each 2-CPU container is CPU-bound;
these are equal-quota containers on one host, not physical machines.

\begin{table*}[t]
\centering
\caption{Two-rank ring in the quota lab at $H=32$ (one seed): bytes sent per
rank and wall time, against one container}
\label{tab:fastsync}
\small
\begin{tabularx}{\textwidth}{X C{2.0cm} C{2.4cm} C{2.4cm} C{2.4cm}}
\toprule
\textbf{Model} & \textbf{One device (s)} & \textbf{MiB per rank, exact / fast} & \textbf{Exact (s), LAN / Wi-Fi} & \textbf{Fast (s), LAN / Wi-Fi} \\
\midrule
MLP-G & 27.9 & 81.1 / 10.7 & 20.0 / 41.8 & 15.2 / 18.9 \\
MLP-C & 16.9 & 98.9 / 10.7 & 14.6 / 43.7 & 10.4 / 14.1 \\
Transformer & 26.7 & 21.7 / 2.4 & 22.7 / 35.9 & 19.1 / 22.9 \\
\bottomrule
\end{tabularx}
\end{table*}
% source: docs/benchmarks/diloco.md (quiet rerun); report Table tab:fastsync.

\begin{figure*}[t]
\centering
\includegraphics[width=0.8\textwidth]{\figpath{fig_5.4_sync_modes.png}}
\caption{Fast synchronisation at $H=32$. Left: quota-lab wall time of one
device and a two-rank ring under exact and fast synchronisation on Wi-Fi, one
seed. Right: in-process validation-loss gap to one device, mean of seeds 0--2
and of seeds 0--5 where run; the dashed line is the 1\% bar.}
\label{fig:sync}
\end{figure*}

At 64 steps, a separate lab run with three seeds per cell also stayed within
the bar on both profiles (worst $+0.64\%$), but in-process the transformer
failed on seeds 0--2 ($+2.13\%$), so the two-rank default is 32. DiLoCo's outer
step left the transformer $21.7\%$ above one device even at one
synchronisation per epoch (156 local steps), the longest interval these runs
allow.

\paragraph{Many ranks} In-process runs set $H(n)$ as the largest interval that
passed on every model: 32 at two ranks, 16 at three, and 4 at four; at eight
ranks the transformer failed at every interval down to 2, so no fast interval
is used for five or more ranks, and fast synchronisation was never measured at
five to seven ranks. Of ten methods tested at 4, 8 and 11 ranks, only exact
synchronisation with rank-16 PowerSGD~\cite{vogels2019powersgd} and error
feedback kept every model within 1\% (worst mean $+0.18\%$), with
$3.2\times$ fewer bytes than 8-bit values on the perceptrons but $1.7\times$
more on the transformer. This method is measured but not the default.
% source: docs/benchmarks/diloco.md.

\paragraph{Adapting the ring} Offered ranks of 16, 8 and 2 CPUs (containers with
those quotas on one host), admission took the 16 and the 8, shared the batch
66:34, and ran in 23.6\,s against 25.4\,s with equal batches. This compares
two rings, not a ring with one machine: in the same recorded session, a
three-rank ring of those containers took 36.4\,s against 6.9\,s for the single
machine that the admission rule picks, and in other runs on that host two CPU
ranks were slower than one (18.9 against 2.8\,s; 191.1 against 92.8\,s).
Throttling one rank's container from 8 CPUs to 1 and back moved the shares
from 64:36 to 84:16 and then to 70:30 after it recovered.
% source: raw/C_commit_measurements.md (2fc3914, 83c59aa, 1c76f51, ac19310).

\subsection{Q5: Physical machines and a phone}
\label{sec:eval-phys}

\begin{table*}[t]
\centering
\caption{Set~A: two laptops in different homes over Tailscale. Script time
unless marked ``from submission''. Sessions~1 and~2 ran different builds.}
\label{tab:phys}
\small
\begin{tabularx}{\textwidth}{X L{6.0cm}}
\toprule
\textbf{Run} & \textbf{Result} \\
\midrule
Mixed CPU script (pool, joblib, NumPy, pandas), one laptop, plain Python, 3 runs & 17.4--18.1\,s \\
Same, two laptops, 3 runs & 13.3--14.6\,s; about half the items on each laptop \\
Whole job sent to the other laptop & 9.1\,s; 11.7\,s from submission \\
Small training model, one GPU & 4.2\,s; validation loss 0.0794 \\
Same, \texttt{-{}-ddp 2}, exact (session 1) & 37.1\,s; loss 0.0797 \\
Same, \texttt{-{}-ddp 2}, fast (session 1) & 7.4\,s; loss 0.0796 \\
Same, \texttt{-{}-ddp 2}, automatic (session 1) & 32.9\,s \\
Same, \texttt{-{}-ddp 2}, fast (session 2) & 8.5\,s; loss 0.0797 \\
Same, \texttt{-{}-ddp 2}, automatic (session 2) & 14.6\,s; loss 0.0796 \\
Combined CPU+GPU script, from submission: first run / repeat & 506.0\,s (487.4\,s building the environment, mostly CUDA PyTorch) / 17.7\,s \\
\bottomrule
\end{tabularx}
\end{table*}
% source: docs/benchmarks/real-hardware.md (exports with run ids).

\paragraph{Set A} Table~\ref{tab:phys} lists the two-laptop runs. The mixed CPU
script ran faster on two laptops than on one (13.3--14.6\,s against
17.4--18.1\,s), with about half its items on the second laptop. The small
training model did not: one GPU finished first (4.2\,s) against every
two-laptop configuration. Fast synchronisation was $5.0\times$ faster than
exact in session~1 (7.4 against 37.1\,s) at about the same loss. Automatic
synchronisation was slow in session~1 (32.9\,s): that build could not switch a
ring fixed by \texttt{-{}-ddp 2} to fast mode. The build of session~2 can; it
synchronised exactly for 19 steps until its first check at the twentieth
synchronisation, then switched (25 synchronisations in all) and took 14.6\,s,
against 8.5\,s for fast synchronisation, whose exact warm-up is 15 of the 147
steps. Exact synchronisation was not repeated in session~2.

\paragraph{Set B} A parallel map took 4.9--5.4\,s on the Core~i9 machine with
the Ryzen laptop as a peer against 5.8--6.4\,s alone, with 116 of 256 items on
the peer; a job of about 5\,s is dominated by setup and serialisation, so the
gain is short of what the core counts allow. Training over the 54\,Mbit/s home
link took 54.9\,s on one machine, 69.4\,s on two ranks with 16-bit gradients
and 56.1\,s with 8-bit gradients (final loss 0.0875, 0.0875, 0.0877): exact
training over that link was not faster than one machine. With a public cloud
host as introducer, the introducer carried 0.2\,MB against about 30\,MB between
the peers, and with the introducer process suspended a known peer was still
reached in 0.03\,s. These figures were each taken a few times by hand and
recorded in the commits that they justified; they were not repeated for this
paper.
% source: docs/benchmarks/raw/C_commit_measurements.md (fca0107, 364b3d9, 42c2426).

\paragraph{A phone} With the S20~FE added to Set~A, a laptop's CPU script spread
over both laptops and the phone gave the phone 8 of 192 items, with an
unchanged checksum, after the phone had built the job's aarch64 environment
from the binary cache in 27\,s. A script training 160 small torch models gave
the phone 6 of them, all of which converged; with 32 models the laptops
finished before the phone's worker had loaded torch, and the phone took none.
Run on the phone alone, the CPU script without its pandas step took 62.4, 62.8
and 66.0\,s under Nix's Python in \texttt{proot}, and 72.8, 73.1 and 80.8\,s
under Termux's own Python. The phone therefore contributes a small share of the
work; we did not measure whether adding it shortened any run.
% source: docs/benchmarks/real-hardware.md "Phone" (commit 950539a).

\paragraph{Across two NATs} On the two-site testbed, every command of the
manual test plan --- scripts, parallel maps, remote runs and both training
synchronisation modes --- ran across the sites and passed its checks, and the
cut-link drill of Table~\ref{tab:drills} recovered the result.
% source: scripts/lab-net-e6.sh; rehearsal.md "Two-site testbed".

% ============================================================
\section{Discussion}
\label{sec:discussion}

\paragraph{Measure, then decline} The recurring result is that distribution
helps only for work that is large relative to its setup and transfer: the
mixed CPU script and the CPU-bound parallel maps gained, while short maps,
small models on a GPU, and exact training over a home link did not. Several
mechanisms exist in their current form because a first version measured slower
than doing nothing: patching libraries at startup ($20\times$ startup), dividing
work by headcount ($1.6\times$ slower than local), returning every rank's
gradient, and synchronous closure broadcast (a 600\,s stall). The system
therefore reports, at the end of each run, how its time was split and where
items ran, and it declines when its models predict a loss.

\paragraph{Models versus measurements} The makespan model guarantees only that
its \emph{prediction} is no worse than local; on the 8:2:1 lab the measured
split was still $1.3\times$ slower than local. The cost model's constants are
hand-set. A model that learns per-machine overheads, and a sensitivity study of
the placement weights of Eq.~\eqref{eq:score}, are open work.

\paragraph{What generalises} Three ideas are not specific to Python: diffing
content-addressed environments by store path between peers; embedding the
introducer in every member so that a small group needs no coordination server;
and choosing a synchronisation interval per ring size from accuracy
measurements rather than from bandwidth alone.

% ============================================================
\section{Threats to validity and limitations}
\label{sec:threats}

\paragraph{External validity} The physical-machine results cover two or three
machines and one phone, mostly single runs or ranges over two or three runs;
one single-sample figure in the project's history was found to be wrong by a
factor of nine on repetition, which is why the repeatable benchmark scripts are
preferred wherever they exist. No wall-time comparison with Ray, Dask, Lithops
or \texttt{torchrun} was run, so Table~\ref{tab:cmp} is qualitative. The labs
share one host: the plain lab cannot show speedup, and the quota lab emulates
quotas and links but not per-core speed, wireless behaviour, campus firewalls
or two physical GPUs. NAT traversal was not exercised between physical homes
(Set~A ran over Tailscale), and the birthday-collision success rate of
Eq.~\eqref{eq:birthday} is computed, not measured.

\paragraph{Internal validity} Some emulated sweeps ran on a loaded host; their
growth fits use CPU time, but their wall times (Table~\ref{tab:scaling}) are
inflated unevenly. The 1\% accuracy bar is a mean over seeds on three small
models, not a guarantee for every seed or for larger models.

\paragraph{System limitations} Only Python on Linux is supported, with NVIDIA
GPUs in practice. The submitter of a parallel script and the machine of a
training ring's rank~0 are single points of failure. Two machines that both
lack a reachable introducer cannot meet without an overlay network, because
there is no relay by design. Phones run chunks only, under a weaker sandbox.
The security gaps of Section~\ref{sec:security} restrict use to mutually
trusting machines.

% ============================================================
\section{Conclusion and future work}
\label{sec:conclusion}

PipedPeer runs unmodified Python on other machines without a configured
cluster: it ships environments as closures diffed by store path, substitutes
the parallel primitives a script already uses, divides work by measured speed,
and joins machines across networks over direct authenticated links with no
relay and no privileged node. On the machines we had, it made CPU-bound work
faster on two machines than on one and kept training accuracy within 1\% at a
fraction of the bytes, and it was slower than one machine for short jobs and
for small models on a GPU. Future work is to bind signing keys to
authenticated identities and require TLS; to let a parallel script and a
training ring survive their submitter and rank~0; to run whole jobs on phones;
to add locality- and topology-aware placement; and to run a controlled,
repeated evaluation on more physical machines against Ray, Dask and Lithops.

% ============================================================
\section*{CRediT authorship contribution statement}
\textbf{Yateen Vaviya:} Conceptualization, Methodology, Software, Validation,
Investigation, Formal analysis, Visualization, Writing -- original draft,
Writing -- review \& editing.
\textbf{Vinayak Yadav:} Conceptualization, Methodology, Software, Validation,
Investigation, Formal analysis, Visualization, Writing -- original draft,
Writing -- review \& editing.
\textbf{Sahil Gupta:} Conceptualization, Methodology, Software, Validation,
Investigation, Formal analysis, Visualization, Writing -- original draft,
Writing -- review \& editing.
\textbf{Rupali Sawant:} Supervision, Writing -- review \& editing.

\section*{Declaration of competing interest}
The authors declare that they have no known competing financial interests or
personal relationships that could have appeared to influence the work reported
in this paper.

\section*{Funding}
This research did not receive any specific grant from funding agencies in the
public, commercial, or not-for-profit sectors.

\section*{Declaration of generative AI and AI-assisted technologies in the manuscript preparation process}
During the preparation of this work the authors used Claude (Anthropic) in
order to assist with drafting and editing the manuscript text and with
developing parts of the software. After using this tool, the authors reviewed
and edited the content as needed and take full responsibility for the content
of the published article.

\section*{Data availability}
The source code, benchmark scripts and every measurement record cited in this
paper are available in the project repository~\cite{pipedpeerrepo}, under
\texttt{docs/benchmarks}.

% ============================================================
\begin{thebibliography}{99}

\bibitem{zaharia2010spark} M.~Zaharia, M.~Chowdhury, M.~J.~Franklin, S.~Shenker, I.~Stoica, Spark: Cluster computing with working sets, in: Proc. 2nd USENIX Workshop on Hot Topics in Cloud Computing (HotCloud), 2010.

\bibitem{moritz2018ray} P.~Moritz, R.~Nishihara, S.~Wang, A.~Tumanov, R.~Liaw, E.~Liang, M.~Elibol, Z.~Yang, W.~Paul, M.~I.~Jordan, I.~Stoica, Ray: A distributed framework for emerging AI applications, in: Proc. 13th USENIX Symp. on Operating Systems Design and Implementation (OSDI), 2018, pp. 561--577.

\bibitem{rocklin2015dask} M.~Rocklin, Dask: Parallel computation with blocked algorithms and task scheduling, in: Proc. 14th Python in Science Conf. (SciPy), 2015, pp. 126--132.

\bibitem{petersohn2020modin} D.~Petersohn, S.~Macke, D.~Xin, W.~Ma, D.~Lee, X.~Mo, J.~E.~Gonzalez, J.~M.~Hellerstein, A.~D.~Joseph, A.~Parameswaran, Towards scalable dataframe systems, Proc. VLDB Endow. 13~(12) (2020) 2033--2046.

\bibitem{burns2016borg} B.~Burns, B.~Grant, D.~Oppenheimer, E.~Brewer, J.~Wilkes, Borg, Omega, and Kubernetes, ACM Queue 14~(1) (2016) 70--93.

\bibitem{litzkow1988condor} M.~J.~Litzkow, M.~Livny, M.~W.~Mutka, Condor -- a hunter of idle workstations, in: Proc. 8th Int. Conf. on Distributed Computing Systems (ICDCS), 1988, pp. 104--111.

\bibitem{zhi2020fiber} J.~Zhi, R.~Wang, J.~Clune, K.~O.~Stanley, Fiber: A platform for efficient development and distributed training for reinforcement learning and population-based methods, arXiv:2003.11164 (2020).

\bibitem{jonas2017pywren} E.~Jonas, Q.~Pu, S.~Venkataraman, I.~Stoica, B.~Recht, Occupy the cloud: Distributed computing for the 99\%, in: Proc. ACM Symp. on Cloud Computing (SoCC), 2017, pp. 445--451.

\bibitem{sampe2021lithops} J.~Samp\'e, P.~Garc\'ia-L\'opez, M.~S\'anchez-Artigas, G.~Par\'is, D.~Barcelona-Pons, D.~Arroyo-Pinto, Toward multicloud access transparency in serverless computing, IEEE Software 38~(1) (2021) 68--74.

\bibitem{arjona2023transparent} A.~Arjona, G.~Finol, P.~Garc\'ia-L\'opez, Transparent serverless execution of Python multiprocessing applications, Future Gener. Comput. Syst. 140 (2023) 436--449. \url{https://doi.org/10.1016/j.future.2022.10.038}

\bibitem{babuji2019parsl} Y.~Babuji, A.~Woodard, Z.~Li, D.~S.~Katz, B.~Clifford, R.~Kumar, L.~Lacinski, R.~Chard, J.~M.~Wozniak, I.~Foster, M.~Wilde, K.~Chard, Parsl: Pervasive parallel programming in Python, in: Proc. 28th Int. Symp. on High-Performance Parallel and Distributed Computing (HPDC), 2019, pp. 25--36.

\bibitem{chard2020funcx} R.~Chard, Y.~Babuji, Z.~Li, T.~Skluzacek, A.~Woodard, B.~Blaiszik, I.~Foster, K.~Chard, funcX: A federated function serving fabric for science, in: Proc. 29th Int. Symp. on High-Performance Parallel and Distributed Computing (HPDC), 2020, pp. 65--76.

\bibitem{dolstra2004nix} E.~Dolstra, M.~de~Jonge, E.~Visser, Nix: A safe and policy-free system for software deployment, in: Proc. 18th Large Installation System Administration Conf. (LISA), 2004, pp. 79--92.

\bibitem{dolstra2006thesis} E.~Dolstra, The Purely Functional Software Deployment Model, Ph.D. thesis, Utrecht University, 2006.

\bibitem{courtes2013guix} L.~Court\`es, Functional package management with Guix, in: Proc. European Lisp Symposium (ELS), 2013.

\bibitem{gamblin2015spack} T.~Gamblin, M.~LeGendre, M.~R.~Collette, G.~L.~Lee, A.~Moody, B.~R.~de~Supinski, S.~Futral, The Spack package manager: Bringing order to HPC software chaos, in: Proc. Int. Conf. for High Performance Computing, Networking, Storage and Analysis (SC), 2015.

\bibitem{merkel2014docker} D.~Merkel, Docker: Lightweight Linux containers for consistent development and deployment, Linux Journal 2014~(239) (2014).

\bibitem{kurtzer2017singularity} G.~M.~Kurtzer, V.~Sochat, M.~W.~Bauer, Singularity: Scientific containers for mobility of compute, PLoS ONE 12~(5) (2017) e0177459.

\bibitem{anderson2002seti} D.~P.~Anderson, J.~Cobb, E.~Korpela, M.~Lebofsky, D.~Werthimer, SETI@home: An experiment in public-resource computing, Commun. ACM 45~(11) (2002) 56--61.

\bibitem{anderson2004boinc} D.~P.~Anderson, BOINC: A system for public-resource computing and storage, in: Proc. 5th IEEE/ACM Int. Workshop on Grid Computing, 2004, pp. 4--10.

\bibitem{cuervo2010maui} E.~Cuervo, A.~Balasubramanian, D.~Cho, A.~Wolman, S.~Saroiu, R.~Chandra, P.~Bahl, MAUI: Making smartphones last longer with code offload, in: Proc. 8th Int. Conf. on Mobile Systems, Applications, and Services (MobiSys), 2010, pp. 49--62.

\bibitem{chun2011clonecloud} B.-G.~Chun, S.~Ihm, P.~Maniatis, M.~Naik, A.~Patti, CloneCloud: Elastic execution between mobile device and cloud, in: Proc. 6th European Conf. on Computer Systems (EuroSys), 2011, pp. 301--314.

\bibitem{arslan2012cwc} M.~Y.~Arslan, I.~Singh, S.~Singh, H.~V.~Madhyastha, K.~Sundaresan, S.~V.~Krishnamurthy, Computing while charging: Building a distributed computing infrastructure using smartphones, in: Proc. 8th Int. Conf. on Emerging Networking Experiments and Technologies (CoNEXT), 2012, pp. 193--204.

\bibitem{fernando2012honeybee} N.~Fernando, S.~W.~Loke, W.~Rahayu, Honeybee: A programming framework for mobile crowd computing, in: Proc. 9th Int. Conf. on Mobile and Ubiquitous Systems (MobiQuitous), 2012, pp. 224--236.

\bibitem{rfc8489} M.~Petit-Huguenin, G.~Salgueiro, J.~Rosenberg, D.~Wing, R.~Mahy, P.~Matthews, Session Traversal Utilities for NAT (STUN), RFC 8489, IETF, 2020.

\bibitem{rfc6887} D.~Wing, S.~Cheshire, M.~Boucadair, R.~Penno, P.~Selkirk, Port Control Protocol (PCP), RFC 6887, IETF, 2013.

\bibitem{rfc6886} S.~Cheshire, M.~Krochmal, NAT Port Mapping Protocol (NAT-PMP), RFC 6886, IETF, 2013.

\bibitem{upnpigd} UPnP Forum, Internet Gateway Device (IGD) V2.0, Standardized DCP, 2010.

\bibitem{rfc8445} A.~Keranen, C.~Holmberg, J.~Rosenberg, Interactive Connectivity Establishment (ICE): A protocol for network address translator (NAT) traversal, RFC 8445, IETF, 2018.

\bibitem{ford2005p2p} B.~Ford, P.~Srisuresh, D.~Kegel, Peer-to-peer communication across network address translators, in: Proc. USENIX Annual Technical Conf., 2005, pp. 179--192.

\bibitem{anderson2020nat} D.~Anderson, How NAT traversal works, Tailscale blog, 2020. URL \url{https://tailscale.com/blog/how-nat-traversal-works} (accessed 16 October 2026).

\bibitem{rfc9000} J.~Iyengar, M.~Thomson, QUIC: A UDP-based multiplexed and secure transport, RFC 9000, IETF, 2021.

\bibitem{donenfeld2017wireguard} J.~A.~Donenfeld, WireGuard: Next generation kernel network tunnel, in: Proc. Network and Distributed System Security Symp. (NDSS), 2017.

\bibitem{graham1969bounds} R.~L.~Graham, Bounds on multiprocessing timing anomalies, SIAM J. Appl. Math. 17~(2) (1969) 416--429.

\bibitem{hochbaum1988uniform} D.~S.~Hochbaum, D.~B.~Shmoys, A polynomial approximation scheme for scheduling on uniform processors: Using the dual approximation approach, SIAM J. Comput. 17~(3) (1988) 539--551.

\bibitem{dean2004mapreduce} J.~Dean, S.~Ghemawat, MapReduce: Simplified data processing on large clusters, in: Proc. 6th USENIX Symp. on Operating Systems Design and Implementation (OSDI), 2004, pp. 137--150.

\bibitem{zaharia2008late} M.~Zaharia, A.~Konwinski, A.~D.~Joseph, R.~Katz, I.~Stoica, Improving MapReduce performance in heterogeneous environments, in: Proc. 8th USENIX Symp. on Operating Systems Design and Implementation (OSDI), 2008, pp. 29--42.

\bibitem{deng2023} Y.~Deng, Z.~Chen, X.~Chen, Y.~Fang, Task offloading in multi-hop relay-aided multi-access edge computing, IEEE Trans. Veh. Technol. 72~(1) (2023) 1372--1376.

\bibitem{yong2023} D.~Yong, R.~Liu, X.~Jia, Y.~Gu, Joint optimization of multi-user partial offloading strategy and resource allocation strategy in D2D-enabled MEC, Sensors 23~(5) (2023) 2565.

\bibitem{he2024} H.~He, X.~Yang, X.~Mi, H.~Shen, X.~Liao, Multi-agent deep reinforcement learning based dynamic task offloading in a device-to-device mobile-edge computing network to minimize average task delay with deadline constraints, Sensors 24~(16) (2024) 5141.

\bibitem{zhao2023} Z.~Zhao, J.~Perazzone, G.~Verma, S.~Segarra, Congestion-aware distributed task offloading in wireless multi-hop networks using graph neural networks, in: Proc. IEEE Int. Conf. on Acoustics, Speech and Signal Processing (ICASSP), 2024, pp. 8951--8955.

\bibitem{li2020pytorchddp} S.~Li, Y.~Zhao, R.~Varma, O.~Salpekar, P.~Noordhuis, T.~Li, A.~Paszke, J.~Smith, B.~Vaughan, P.~Damania, S.~Chintala, PyTorch distributed: Experiences on accelerating data parallel training, Proc. VLDB Endow. 13~(12) (2020) 3005--3018.

\bibitem{sergeev2018horovod} A.~Sergeev, M.~Del~Balso, Horovod: Fast and easy distributed deep learning in TensorFlow, arXiv:1802.05799 (2018).

\bibitem{dettmers2016eightbit} T.~Dettmers, 8-bit approximations for parallelism in deep learning, in: Proc. Int. Conf. on Learning Representations (ICLR), 2016.

\bibitem{seide2014onebit} F.~Seide, H.~Fu, J.~Droppo, G.~Li, D.~Yu, 1-bit stochastic gradient descent and its application to data-parallel distributed training of speech DNNs, in: Proc. Interspeech, 2014, pp. 1058--1062.

\bibitem{vogels2019powersgd} T.~Vogels, S.~P.~Karimireddy, M.~Jaggi, PowerSGD: Practical low-rank gradient compression for distributed optimization, in: Advances in Neural Information Processing Systems (NeurIPS), 2019.

\bibitem{stich2019localsgd} S.~U.~Stich, Local SGD converges fast and communicates little, in: Proc. Int. Conf. on Learning Representations (ICLR), 2019.

\bibitem{lin2020localsgd} T.~Lin, S.~U.~Stich, K.~K.~Patel, M.~Jaggi, Don't use large mini-batches, use local SGD, in: Proc. Int. Conf. on Learning Representations (ICLR), 2020.

\bibitem{mcmahan2017fedavg} B.~McMahan, E.~Moore, D.~Ramage, S.~Hampson, B.~A.~y~Arcas, Communication-efficient learning of deep networks from decentralized data, in: Proc. 20th Int. Conf. on Artificial Intelligence and Statistics (AISTATS), 2017, pp. 1273--1282.

\bibitem{douillard2023diloco} A.~Douillard, Q.~Feng, A.~A.~Rusu, R.~Chhaparia, Y.~Donchev, A.~Kuncoro, M.~Ranzato, A.~Szlam, J.~Shen, DiLoCo: Distributed low-communication training of language models, arXiv:2311.08105 (2023).

\bibitem{ryabinin2020learningathome} M.~Ryabinin, A.~Gusev, Towards crowdsourced training of large neural networks using decentralized mixture-of-experts, in: Advances in Neural Information Processing Systems (NeurIPS), 2020.

\bibitem{diskin2021dedloc} M.~Diskin, A.~Bukhtiyarov, M.~Ryabinin, L.~Saulnier, Q.~Lhoest, A.~Sinitsin, D.~Popov, D.~Pyrkin, M.~Kashirin, A.~Borzunov, A.~Villanova~del~Moral, D.~Mazur, I.~Kobelev, Y.~Jernite, T.~Wolf, G.~Pekhimenko, Distributed deep learning in open collaborations, in: Advances in Neural Information Processing Systems (NeurIPS), 2021.

\bibitem{borzunov2023petals} A.~Borzunov, D.~Baranchuk, T.~Dettmers, M.~Ryabinin, Y.~Belkada, A.~Chumachenko, P.~Samygin, C.~Raffel, Petals: Collaborative inference and fine-tuning of large models, in: Proc. 61st Annual Meeting of the Association for Computational Linguistics: System Demonstrations (ACL), 2023, pp. 558--568.

\bibitem{kth2025thesis} J.~Spenger, Programming Models for Failure-Transparent Distributed Systems, Ph.D. thesis, KTH Royal Institute of Technology, Stockholm, 2025.

\bibitem{thaler1998rendezvous} D.~G.~Thaler, C.~V.~Ravishankar, Using name-based mappings to increase hit rates, IEEE/ACM Trans. Netw. 6~(1) (1998) 1--14.

\bibitem{bernstein2012ed25519} D.~J.~Bernstein, N.~Duif, T.~Lange, P.~Schwabe, B.-Y.~Yang, High-speed high-security signatures, J. Cryptogr. Eng. 2~(2) (2012) 77--89.

\bibitem{rfc5705} E.~Rescorla, Keying material exporters for Transport Layer Security (TLS), RFC 5705, IETF, 2010.

\bibitem{rfc2104} H.~Krawczyk, M.~Bellare, R.~Canetti, HMAC: Keyed-hashing for message authentication, RFC 2104, IETF, 1997.

\bibitem{rfc8446} E.~Rescorla, The Transport Layer Security (TLS) protocol version 1.3, RFC 8446, IETF, 2018.

\bibitem{cloudpickle} cloudpickle: Extended pickling support for Python objects, 2015. URL \url{https://github.com/cloudpipe/cloudpickle} (accessed 16 October 2026).

\bibitem{crun} crun: A fast and lightweight fully featured OCI runtime and C library for running containers, 2017. URL \url{https://github.com/containers/crun} (accessed 16 October 2026).

\bibitem{termux} Termux: Android terminal emulator and Linux environment, 2015. URL \url{https://termux.dev} (accessed 16 October 2026).

\bibitem{proot} PRoot: chroot, mount --bind, and binfmt\_misc without privilege/setup for Linux, 2010. URL \url{https://proot-me.github.io} (accessed 16 October 2026).

\bibitem{jia2011mutation} Y.~Jia, M.~Harman, An analysis and survey of the development of mutation testing, IEEE Trans. Softw. Eng. 37~(5) (2011) 649--678.

\bibitem{pipedpeerrepo} Y.~Vaviya, V.~Yadav, S.~Gupta, PipedPeer: source code and measurement records. URL \url{https://github.com/vinayakyadav2709/Pipedpeer} (accessed 16 October 2026).

\end{thebibliography}

\end{document}
